\scanpage{103}
\mspage{52}
\uline{Définition 4} : On dit que $f:\deleted{\unsure{S}}X\longrightarrow Y$ dans $\mathrm{Cat}$ est une $W$-fibration, si elle satisfait les conditions \inserted{\unsure{équiv.|préc.} a) à $c_2$} de la prop. 9.

\uline{Corollaire 1} Si $f$ est \inserted{$W$-}propre ou \inserted{$W$-}lisse, pour que $f$ soit une $W$-fibration, il faut et \unsure{s|s.} que $f$ soit une $W$-pseudofibration, i.e. que pour $S'=\Delta_1\longrightarrow S$, l\deleted{\unsure{'ori} \unsure{des}}'inclusion des deux fibres
\diagram{\begin{tikzcd}[ampersand replacement=\&]
X'_0\arrow[r,"i_0"] \& X' \& X'_1\arrow[l,"i_1"']
\end{tikzcd}}
de $X'=X\times_S S'$ soient dans $W$.

\uline{Corollaire 2} Supposons $f$ Cat-fibrant. Alors $f$ est une $W$-fibration sss pour \unsure{ts} flèche $u$ \inserted{$:s\longrightarrow s'$} dans $S$, $u^*:X_{s'}\longrightarrow X_s$ est dans $W$. Énoncé dual si $f$ est Cat-cofibrant.

\uline{Corollaire 3} Supposons $f$ $W$-prépropre et $W$-prélisse (par exemple $f$ $W$-parfait). Alors $f$ est une $W$-fibration.
% margin note written sideways
\margin{\uline{Prop.} 10 Soit $f:X\to Y$ une flèche dans \deleted{\unsure{$\mathrm{Cat}$}} $\underline{\mathrm{Fib}}_W S^*$ (sous-cat. pleine de $\mathrm{Cat}_{/S}$ formée des $W$-fibrés sur $S$). Si $\forall s\in S$, $f_s:X_s\to Y_s$ est dans $W$, alors $f\in W$. \deleted{\unsure{Il suffit de}} Considérons \diagram{\begin{tikzcd}[ampersand replacement=\&,column sep=small,row sep=small] X\arrow[r,"f"]\arrow[dr,"{i_X}"']\&Y\arrow[dr,"{i_Y}"]\&{}\\ \&\Phi_S X\arrow[r,"{\Phi_S(f)}"']\&\Phi_S Y\end{tikzcd}}, si $i_X,i_Y,f\in W$, alors $\Phi_S(f)\in W$ \deleted{$\Rightarrow$ \unsure{$f$ est de}} $\Rightarrow f\in W$. Or le fait que $X,$\deleted{\unsure{$Y$}} $W$-fibrés \deleted{\unsure{, sur}} $S$ implique que $i_X\in W_S^f$ (il suffit $=$ que $X$ soit prépropre sur \unsure{$S$}) idem pour $Y,i_Y$. Donc l'hyp. $f\in W_S^f$ implique $\underline{\Phi}_S(f)\in W_S^f$. Donc $\Phi_S(f)\in W$ par prop. 5 (page 29).}
% margin note written sideways
\margin{$^*$ $X$ et $Y$ tous deux prépropres, ou tous deux prélisses \unsure{suffit|suffisent}}

\deleted{\uline{Théorème 4} Soit $W\subset\mathrm{Fl}(\mathrm{Cat})$ satisfaisant $W(1,2,3)$. Considérons les conditions suivantes sur $W$ :}
\deleted{a) Toute $W$-fibration est ultraforte, i.e. \inserted{si $f:X\to S$ dans $\mathrm{Cat}$,} les conditions équiv. de la prop. 9 sur $f$ impliquent que pour $S''\overset{u}{\unsure{\relbar|\longrightarrow}}S'$ dans $\mathrm{Cat}_{/S}$, si $u\in W$, alors $\inserted{u_X:}X''\longrightarrow X'$ est $\unsure{\in}W$.}
\deleted{a$'$) \deleted{Toute \unsure{$W$-fib}} Tout $f$ \inserted{$:X\to S$,} qui est $W_{\unsure{\infty}}$-parfait (donc a fortiori $W_{\deleted{\unsure{\infty}}}$-parfait, donc une $W$-fibration), est une $W$-fibration ultraforte.}
\deleted{b) Pour tte $S$ dans $\mathrm{Cat}$ et toute flèche $f:X\to Y$ dans $\mathrm{Cat}_{/\deleted{\unsure{S}}S}$, si $X,Y$ sont $W$-fibrés sur $S$, et $f\in W$,}
% archive sheet number: 103
